\scanpage{6}
\mspage{\handwritten{9}}
équivalentes.

L'équivalence de (ii) (iii) (iv) (en fait, pour tout $a\handwritten{\in}W(k)$ choisi) et le fait que ces conditions impliquent toutes les autres, \deleted{\unsure{e}}est contenu dans 1.3. En vertu de 1.2. le fait que $N$ est lisse sur $k$ et (2) une égalité équivaut à $\dim X=\dim V$, ce qui montre l'équivalence de (i bis) et (i ter), qui impliquent évidemment (i), reste à prouver que (i) implique (i bis), et , lorsque $k$ est algébriquement clos, que (i) implique (ii). Quitte à passer à la clôture algébrique de $k$, on peut supposer $k$ algébriquement clos , et on est réduit à prouver alors $\text{(i)}\handwritten{\Rightarrow}\text{(ii)}$. Or comme $\dim X\handwritten{\geq}\dim V$ en vertu de 1.2. , on voit que (i) implique que $\deleted{\phi}\handwritten{\psi}:X\handwritten{\to}V$ est dominant , et comme $V$ est génériquement réduit, \deleted{\unsure{il est génériquement régulier, il existe donc un ouvert}}\inserted{il e\handwritten{x}iste un ouvert non vide de $X$ qui est régulier et à} \unsure{fortiori|fortiooi} normal, donc il existe un \deleted{\unsure{$x\ X$ tel}} $x\handwritten{\in}X(k)$ tel que $\deleted{\phi}\handwritten{\psi}$ soit non ramifié en $x$ et que $\deleted{\phi}\handwritten{\psi}(x)$ soit un point \deleted{normal de $V$, \unsure{alors|Alors}} régulier et à fortiori normal de $V$. On aura $x=\deleted{\unsure{a|\phi}}(g^{\handwritten{\unsure{\circ|\bullet}}},a)$, et quitte à remplacer $x$ par $(e^{\handwritten{\circ}},g^{-1}.a)$, on peut supposer $a\handwritten{\in}W(k)$. Alors en vertu de(1.3. $\text{(iv)}\handwritten{\Rightarrow}\text{(i)}$ ) il s'ensuit que $a$ satifait à la condition 1.5. (ii), ce qui achève la démonstration.

\uline{Corollaire} 1.6. Soient $G,V,W$ comme dans 1.\deleted{\unsure{5}}\handwritten{3}. \deleted{\unsure{Il}}\inserted{\handwritten{$V$ normal, $k$} alg. clos.} Pour que $\handwritten{\psi}:X\handwritten{\to}V$ soit un morphisme \uline{birationnel}, il faut et il suffit que les conditions équivalentes de 1.5. soient satisfaites, et que de plus il existe un ouvert non vide $U$ de $V$ tel que tout $a\handwritten{\in}U(k)$ soit contenu dans un unique conjugué de $W$ par un élé\handwritten{ment} de $G(k)$. De plus, pour un $a\handwritten{\in}V(k)$, les conditions suivantes sont alors équivalentes:

(i) Il existe un unique conjugué de $W$ contenant $a$.

(ii) L'ensemble des conjugués de $W$ contenant $a$ est fini non vide.

(iii) Le schéma $\handwritten{\psi}^{-1}(a)$ \handwritten{``}des conjugués de $W$ qui contiennent $a$\handwritten{''} contient un point isolé.
% archive sheet number: 5
